\chapter{Sequential decisions and dynamic programming}
\label{sec:p17-sequential-decisions}

A sequential decision changes both the current payoff and the
distribution of later opportunities. The method below takes the
transition law, information available before each action, and an
additive payoff criterion as inputs. It returns an optimal policy
and a value, with computable error bounds when iteration is used.
The source interfaces are Bellman's bounded transformation equation
\cite[Theorem~1, pp.~716--718]{Bellman1952DynamicProgramming}
and Blackwell's discounted model and operator construction
\cite[Sections~2--3 and Theorems~3--8, pp.~227--234]{Blackwell1965Discounted}.
Bellman's theorem is proved at its stated bounded-function scope.
The stochastic decision results here use finite states and finite
actions; Blackwell's broader Borel-space results have additional scope.
The remaining six original Bellman constructions, with their full
resource, stopping and search domains, are supplied in
\cref{sec:p17-bellman-original-constructions}.

\section{What is observed, chosen and received}

\begin{hbdefinition}{def:p17-sequential-model}{A finite controlled Markov model}
Let \(S\) be a nonempty finite state set. At time
\(t=0,1,\ldots\), state \(s\) admits the nonempty finite action
set \(A_t(s)\). A finite common set \(A\) contains these actions.
The inputs are transition probabilities and real rewards
\[
 P_t(s'\mid s,a)\geq0,\qquad
 \sum_{s'\in S}P_t(s'\mid s,a)=1,\qquad
 r_t(s,a,s')\in\mathbb R
 \quad(a\in A_t(s)).
\]
At time \(t\), the decision maker first observes \(X_t\),
then chooses \(A_t\), then the next state \(X_{t+1}\) is
drawn from \(P_t(\cdot\mid X_t,A_t)\), and the reward
\(r_t(X_t,A_t,X_{t+1})\) is received.
The symbol \(A_t\) for the chosen random action is distinguished
by its lack of the state argument from the feasible set \(A_t(s)\).
The expected one-step reward is
\[
                 \bar r_t(s,a)=
                     \sum_{s'}P_t(s'\mid s,a)r_t(s,a,s').
\]
The controlled Markov assumption specifies this next-state law
even conditional on the whole available history and current action.
It also specifies that the displayed state/action/next-state
arguments contain all information relevant to the current reward.

A history before action \(t\) is
\(h_t=(s_0,a_0,s_1,\ldots,a_{t-1},s_t)\), with every past
action feasible at its recorded state and time. A
\emph{randomized history-dependent policy} \(\pi\) specifies
probabilities \(\pi_t(a\mid h_t)\) on \(A_t(s_t)\) at every
such history. A deterministic policy puts mass one on a selected
action. A Markov policy depends only on \(t,s_t\).
In a time-homogeneous model, \(A_t(s)=A(s)\),
\(P_t=P\) and \(r_t=r\); a deterministic stationary policy
is a single map \(f\) with \(f(s)\in A(s)\).
Randomization uses no signal about the environment beyond
the stated history.

An initial distribution \(\rho\in\Delta(S)\) may be supplied;
the value from initial state \(s\) uses \(\rho=\delta_s\).
The expected-utility interpretation uses the finite lotteries
and signalwise comparisons in
\cref{thm:p17-lottery-reduction,thm:p17-information-policy-value}.
Additivity across time and the chosen discount are further
modeling assumptions. The finite-prize representation
of \cref{thm:p17-finite-vnm} alone does not impose them.
\index{Markov decision process}\index{policy!history-dependent}
\index{dynamic programming}
\end{hbdefinition}

\begin{hbtheorem}{thm:p17-sequential-policy-law}{proved=annex:p17.sequential.policy.law}{The probability law of every allowed policy}
For each initial law and policy there is a unique path probability
law with cylinder masses
\[
 \rho(s_0)\prod_{t=0}^{N-1}
          \pi_t(a_t\mid h_t)P_t(s_{t+1}\mid s_t,a_t).
\]
All finite-horizon expectations are therefore finite sums.
The same specified kernels define a continuation from any
history, even a history having zero probability from the chosen
initial law. Statements about such continuations use this
specified law rather than a division by a zero history probability.
For any fixed initial law, a randomized Markov policy can reproduce
the state/action marginals of a history-dependent policy at every
time. This latter policy may depend on the initial law; a single
simultaneously optimal policy for all initial states is established
separately below.
\end{hbtheorem}

\section{A finite number of decisions}

\begin{hbdefinition}{def:p17-sequential-finite-return}{Horizon, terminal payoff and continuation value}
Fix an integer \(N\geq0\), terminal reward \(g:S\to\mathbb R\),
and \(0\leq\gamma\leq1\). There are decisions at
\(t=0,\ldots,N-1\), followed by \(g(X_N)\).
Only those \(N\) transition, reward and feasible-action tables
are required; an arbitrary feasible extension beyond \(N\)
has no effect on this return.
For a continuation from \(h_t\), let
\[
 J_{t,N}^{\pi}(h_t)=
 E_{h_t}^{\pi}\!\left[
   \sum_{k=t}^{N-1}\gamma^{k-t}
           r_k(X_k,A_k,X_{k+1})+\gamma^{N-t}g(X_N)\right].
\]
Use \(\gamma^0=1\), including \(\gamma=0\), and an empty sum
equals zero. At \(t=N\), this return is \(g(s_N)\).
The finite-horizon optimum is the supremum over all specified
history-dependent randomized continuation policies.
\end{hbdefinition}

\begin{hbtheorem}{thm:p17-sequential-backward}{proved=annex:p17.sequential.backward}{Backward recursion and full policy-class optimality}
Set \(v_N=g\) and compute, for \(t=N-1,\ldots,0\),
\[
 v_t(s)=\max_{a\in A_t(s)}
       \sum_{s'}P_t(s'\mid s,a)
                       [\,r_t(s,a,s')+\gamma v_{t+1}(s')\,].
\]
Every maximum is attained. Select a maximizing action
\(f_t(s)\) at each state and time, using a fixed order to resolve
ties. Then for every history ending in \(s\) at time \(t\),
\[
                 J_{t,N}^{\pi}(h_t)\leq v_t(s)
                         =J_{t,N}^{f}(h_t)
                    \quad\hbox{for every }\pi.
\]
Thus a deterministic Markov policy is simultaneously optimal
from every history. For initial law \(\rho\), its value is
\(\sum_s\rho(s)v_0(s)\).
When \(N=0\), the answer is \(g\) and there are no actions.
With rational input tables, every value and selected action can
be obtained by finite exact rational arithmetic and comparison.
\end{hbtheorem}

\section{Discounted infinite time}

\begin{hbdefinition}{def:p17-sequential-discounted}{Bounded return and the two operators}
For the discounted stochastic theorems below the model is time homogeneous,
\(0\leq\gamma<1\), and
\[
 R=\max_{\substack{s,s'\in S\\a\in A(s)}}|r(s,a,s')|<\infty .
\]
This boundedness follows from the finite input tables; it is
recorded explicitly because the infinite-series proof uses it.
Define
\[
 J^\pi(s)=E_s^\pi\sum_{t=0}^{\infty}
                      \gamma^t r(X_t,A_t,X_{t+1}),\qquad
 V^*(s)=\sup_\pi J^\pi(s).
\]
For \(v:S\to\mathbb R\), write
\(\|v\|_\infty=\max_s|v(s)|\), and set
\[
\begin{aligned}
 (Tv)(s)&=\max_{a\in A(s)}
                  [\,\bar r(s,a)+\gamma\sum_{s'}P(s'\mid s,a)v(s')\,],\\
 (T_fv)(s)&=\bar r(s,f(s))
                         +\gamma\sum_{s'}P(s'\mid s,f(s))v(s').
\end{aligned}
\]
\(T\) is the optimality operator; \(T_f\) fixes the first action
to the stationary choice \(f\).
\end{hbdefinition}

\begin{hbtheorem}{thm:p17-sequential-tail}{proved=annex:p17.sequential.tail}{Convergent returns and uniform tail bounds}
Every displayed discounted reward series converges absolutely
on every admissible path, and
\[
 |J^\pi(s)|\leq\frac R{1-\gamma},\qquad
 \left|J^\pi(s)-E_s^\pi\sum_{t=0}^{N-1}\gamma^t r_t\right|
                       \leq\frac{R\gamma^N}{1-\gamma}.
\]
Here \(r_t\) abbreviates the realized reward of the homogeneous
model. If two policies have the same decision kernels at all
histories at times \(0,\ldots,N-1\), then
\[
                  \|J^\pi-J^{\pi'}\|_\infty
                          \leq\frac{2R\gamma^N}{1-\gamma}.
\]
The factor \(2\) is necessary for rewards of both signs.
With terminal payoff \(v\), the corresponding \(N\)-stage
return differs from \(J^\pi\) by at most
\(\gamma^N(R/(1-\gamma)+\|v\|_\infty)\).
\end{hbtheorem}

\begin{hblemma}{thm:p17-sequential-contraction}{proved=annex:p17.sequential.contraction}{The full contraction and fixed-point argument}
For any nonempty set \(E\), bounded functions on \(E\)
form a complete space in the norm
\(\|u\|_\infty=\sup_{x\in E}|u(x)|\).
If an operator \(B\) on this space satisfies
\[
                 \|Bu-Bv\|_\infty\leq\gamma\|u-v\|_\infty,
                         \qquad 0\leq\gamma<1,
\]
then it has a unique fixed point \(w\). Its successive iterates
\(u_{k+1}=Bu_k\) satisfy
\[
 \|u_k-w\|_\infty\leq\gamma^k\|u_0-w\|_\infty,\qquad
 \|u_k-w\|_\infty\leq
               \frac{\gamma^k}{1-\gamma}\|Bu_0-u_0\|_\infty.
\]
The same conclusion holds on a nonempty closed invariant subset.
For any specified sigma algebra on \(E\), its bounded measurable
functions form such a closed space whenever the operator preserves them.
Here monotonicity means that pointwise \(u\leq v\) implies
pointwise \(Bu\leq Bv\).
In particular, an operator which is monotone and satisfies
\[
                  B(u+c\mathbf1)=Bu+\gamma c\mathbf1
                           \quad(c\in\mathbb R)
\]
is such a contraction. The finite operators \(T,T_f\) have
these properties.
\end{hblemma}

\begin{hbtheorem}{thm:p17-bellman-bounded-transformations}{proved=annex:p17.bellman.bounded.transformations}{Bellman's original bounded transformation equation}
Let \(E\ne\varnothing\), fix an integer \(m\geq1\),
and let \(S_k:E\to E\), \(k=1,\ldots,m\).
Suppose
\[
       0\leq g_k(p)\leq C,\qquad
       0\leq h_k(p)\leq b<1
                  \quad(p\in E,\ 1\leq k\leq m).
\]
Then the equation
\[
            F(p)=\max_{1\leq k\leq m}
                      [\,g_k(p)+h_k(p)F(S_kp)\,]
\]
has exactly one bounded solution. It is nonnegative, bounded
by \(C/(1-b)\), and obtained by successive approximation
from any bounded initial function.
This proves Bellman's Theorem~1, whose \(E\) is a region
of Euclidean space, without requiring \(E\) to be finite.
When \(E\) is finite, for constant \(h_k=\gamma\) and deterministic transitions
\(p\mapsto S_kp\), this is the optimality operator above
with nonnegative rewards. The stochastic operator replaces
evaluation at a successor by its probability-weighted average.
\end{hbtheorem}

\begin{hbtheorem}{thm:p17-sequential-stationary-optimum}{proved=annex:p17.sequential.stationary.optimum}{A stationary optimum against all history-dependent policies}
The unique fixed point \(w=Tw\) is \(V^*\).
At each state choose a maximizing action \(f^*(s)\) in
\(Tw(s)\). Then
\[
               J^\pi(s)\leq w(s)=J^{f^*}(s)
                       \quad(s\in S,\ \hbox{every }\pi).
\]
Thus \(f^*\) is deterministic, stationary, and optimal from
every state and every initial distribution.
If a vector \(u\) satisfies \(Tu\leq u\), it is an upper
bound on every policy value. For a stationary \(f\),
\[
              f\hbox{ is optimal}
          \quad\Longleftrightarrow\quad TV^f=V^f,
                 \qquad V^f=J^f.
\]
More generally a policy is optimal from all initial states
if and only if its value vector is a fixed point of \(T\).
\end{hbtheorem}

\begin{hbtheorem}{thm:p17-sequential-value-iteration}{proved=annex:p17.sequential.value.iteration}{Value iteration and a policy-loss certificate}
Starting from any vector \(v_0\), iterate \(v_{k+1}=Tv_k\).
For \(v_0=0\),
\[
                   \|v_k-V^*\|_\infty
                        \leq\frac{R\gamma^k}{1-\gamma}.
\]
For any vector \(v\), define its residual
\(\delta(v)=\|Tv-v\|_\infty\). Then
\[
                       \|v-V^*\|_\infty
                             \leq\frac{\delta(v)}{1-\gamma}.
\]
If a stationary policy \(f\) obeys
\(T_fv\geq Tv-\eta\mathbf1\), \(\eta\geq0\), then
\[
 0\leq V^*(s)-V^f(s)
       \leq\frac{2\gamma\|v-V^*\|_\infty+\eta}{1-\gamma}
       \leq
          \frac{2\gamma\,\delta(v)}{(1-\gamma)^2}
                         +\frac{\eta}{1-\gamma}.
\]
An exact greedy choice has \(\eta=0\). These bounds give
explicit stopping criteria for value accuracy and for the
loss of the extracted policy. At \(\gamma=0\), one
application of \(T\) gives the optimal value and a greedy
policy is optimal.
\end{hbtheorem}

\section{Evaluating and improving a policy}

\begin{hbtheorem}{thm:p17-sequential-policy-evaluation}{proved=annex:p17.sequential.policy.evaluation}{The policy linear system and its actual inverse}
For a stationary \(f\), put
\[
 P_f(s,s')=P(s'\mid s,f(s)),\qquad r_f(s)=\bar r(s,f(s)).
\]
Every power of \(P_f\) is row stochastic, and
\[
 (I-\gamma P_f)^{-1}=\sum_{k=0}^{\infty}\gamma^kP_f^k,
 \qquad
 V^f=(I-\gamma P_f)^{-1}r_f.
\]
The inverse exists, has nonnegative entries, and has
maximum absolute row sum \(1/(1-\gamma)\).
Its tail after power \(N-1\) has row sum
\(\gamma^N/(1-\gamma)\).
The value is the unique solution of \(v=T_fv\).
For rational transition, reward and discount data, Gaussian
elimination computes this vector exactly using finitely many
rational operations; each required pivot exists after a row
interchange because the matrix is invertible.
\end{hbtheorem}

\begin{hbtheorem}{thm:p17-sequential-improvement}{proved=annex:p17.sequential.improvement}{Both source improvement routines}
Let \(v=J^\pi\) for any policy, and let \(g\) be stationary.
If \(T_gv\geq v\), then
\[
                        V^g\geq T_gv\geq v.
\]
Here \(T_gJ^\pi\) is the return from using \(g\) for the
first action and then restarting \(\pi\) from the observed
successor state.
In particular, a greedy policy \(g\) for \(V^f\) improves
a stationary \(f\). Writing \(d=T_gV^f-V^f\geq0\),
\[
               V^g-V^f=(I-\gamma P_g)^{-1}d\geq d.
\]
Thus an improvement strict at a state gives a value increase
strict at that same state.

For two stationary policies \(f,g\), define \(h(s)=f(s)\)
where \(V^f(s)\geq V^g(s)\), and \(h(s)=g(s)\) otherwise.
Then
\[
                         V^h\geq\max(V^f,V^g)
\]
coordinatewise. These are, respectively, the Howard and
Eaton--Zadeh improvement statements proved in Blackwell's
Theorem~8, specialized to the present finite model.
\end{hbtheorem}

\begin{hbtheorem}{thm:p17-sequential-policy-iteration}{proved=annex:p17.sequential.policy.iteration}{Finite termination with a specified tie rule}
Start with any deterministic stationary policy \(f\).
Evaluate \(V^f\). At each state retain \(f(s)\) if it
maximizes the action value computed from \(V^f\); otherwise
replace it by the first maximizing action in a fixed action
order. Repeat evaluation and improvement until no action changes.
The final policy is optimal, and there are at most
\(\prod_{s\in S}|A(s)|-1\) strict improvement steps.
This is a terminating exact algorithm for rational data.
For arbitrary unspecified real inputs, the mathematical
maxima and inverses still exist, but an effective exact
comparison procedure for those real numbers is an additional input.
\end{hbtheorem}

\section{A complete maintenance decision}

\begin{hbdefinition}{def:p17-sequential-maintenance}{Two observed machine conditions}
The states are \(G\), good, and \(B\), bad. The following
table gives the reward, which here is independent of the
successor state, and the next-state probabilities.
\[
\begin{array}{c|c|r|cc}
 \text{state}&\text{action}&r&P(G)&P(B)\\ \hline
 G&\text{run}&3&1/2&1/2\\
 G&\text{service}&1&1&0\\
 B&\text{wait}&0&0&1\\
 B&\text{repair}&-2&1&0
\end{array}
\]
The condition is observed before acting each period. Running
earns three now and may leave the machine bad for tomorrow.
Service earns one now and leaves it good. Repair incurs two
now and makes the next period good; it does not also produce
the running reward in that period. These timing choices are
part of the model.
\end{hbdefinition}

\begin{hbtheorem}{thm:p17-sequential-maintenance}{proved=annex:p17.sequential.maintenance}{Backward decisions, exact policy evaluation and an error certificate}
For three decisions, terminal reward zero and \(\gamma=1\),
the backward values and selected policies are
\[
\begin{array}{c|cc|cc}
 t&v_t(G)&v_t(B)&f_t(G)&f_t(B)\\ \hline
 3&0&0&\text{no action}&\text{no action}\\
 2&3&0&\text{run}&\text{wait}\\
 1&9/2&1&\text{run}&\text{repair}\\
 0&23/4&5/2&\text{run}&\text{repair}
\end{array}
\]
Starting in \(G\), the optimal value is \(23/4\).
The decision after a bad transition uses the newly observed
condition; repair is worthwhile with two decisions left
and is wasteful on the last decision.

For discounted infinite time with \(\gamma=3/4\), policy
\((\text{run},\text{wait})\) has values \((24/5,0)\).
One greedy improvement gives
\((\text{run},\text{repair})\), with
\[
                         V^*=(72/11,\ 32/11).
\]
The two omitted action values at this vector are
\(65/11\) for service in \(G\) and \(24/11\) for wait in \(B\),
so the displayed policy is optimal.
Value iteration from zero begins
\[
              v_1=(3,0),\qquad v_2=(33/8,1/4).
\]
Here \(R=3\). After \(k\) iterations its value error is at
most \(12(3/4)^k\), and an exact greedy stationary policy
has loss at most \(72(3/4)^k\).
In particular \(k=23\) certifies loss below \(1/10\);
the direct fixed-point calculation above certifies the
specific displayed policy with zero loss.
\end{hbtheorem}

\section{The boundaries of the method}

\begin{hbtheorem}{thm:p17-sequential-boundaries}{proved=annex:p17.sequential.boundaries}{What fails after the assumptions change}
The following examples have the stated outcomes.

\begin{enumerate}
\item With one state, one action and reward \(1\), discount
\(\gamma=1\) gives infinite total reward and equation \(v=1+v\)
has no finite solution. With reward zero instead, every
constant is a fixed point. With two deterministically alternating
states and alternating rewards \(1,-1\), the undiscounted
series is undefined as a convergent real series.
\item Boundedness is automatic for the finite homogeneous tables.
If the state space is enlarged to \(\mathbb N_0\), with the
only transition \(j\mapsto j+1\), reward \(2^j\) and
\(\gamma=1/2\), the return from zero is infinite.
\item In Blackwell's one-state example with actions
\(a\in\{1,2,\ldots\}\), reward \(1-1/a\), and
\(0\leq\gamma<1\), the supremum is \(1/(1-\gamma)\)
but no history-dependent randomized policy attains it.
Thus bounded rewards alone do not replace finite action attainment.
\item If a fair bit \(H\) is observed at time zero and then
the recorded state at time one is the same symbol for both
bit values, a final action \(a\in\{0,1\}\) rewarded by
\(\mathbf1_{\{a=H\}}\) has history-policy value \(1\).
A policy using only that time-one symbol has value at most
\(1/2\). The recorded state fails the reward sufficiency
assumption. Retaining \(H\) in the state repairs this particular model.
\item With one state and actions \(0,1\), rewards
\(r_t(a)=\mathbf1_{\{a=t\bmod2\}}\), and \(\gamma=1/2\),
a time-dependent Markov policy has value \(2\).
Every stationary randomized policy has value at most \(4/3\).
This example removes time homogeneity.
\item With one state, two identical zero-reward actions and
\(\gamma<1\), switching actions on each tie can continue forever.
The specified retain-on-tie rule prevents this policy-iteration failure.
\end{enumerate}
\end{hbtheorem}

\begin{hblemma}{thm:p17-sequential-source-bridge}{proved=annex:p17.sequential.source.bridge}{The exact source operator and action conventions}
Blackwell's reward is \(r(s,a,s')\), received after the
transition; his discount is \(\beta=\gamma\), with the
first reward undiscounted. With common finite action set
and finite states, his stationary operator is \(T_f\).
The supremum operator generated by an enumeration of all
common actions is \(T\).
State-dependent feasible action sets can be translated into
his common-action convention without changing attainable
state/reward laws or values.

For \(\gamma>0\), two actions at a fixed state define the
same one-step operator on every continuation vector exactly
when they have the same transition row and the same
expected reward \(\bar r(s,a)\). Equality of the realized
reward function at every successor is sufficient but is not
necessary. At \(\gamma=0\), equality of expected reward
alone is equivalent to equality of the operators.
\end{hblemma}

\begin{hbmethod}{def:p17-sequential-computation}{Inputs, computation and returned decision}
Supply the finite state and feasible-action lists, transition rows,
successor-dependent rewards, and observation/action timing.
For a finite horizon also supply \(N,g,\gamma\); apply
\cref{thm:p17-sequential-backward} to return all \(v_t,f_t\)
and the initial-distribution value.
For the homogeneous discounted criterion supply \(0\leq\gamma<1\).
With exact rational data, \cref{thm:p17-sequential-policy-iteration}
returns an exact optimal stationary policy and its evaluated value.
Alternatively specify an error tolerance, use
\cref{thm:p17-sequential-value-iteration}, and return the
computed value, greedy policy, and the corresponding residual
or geometric error certificate.
All numerical input values, including discount and terminal reward,
must have the representation needed for the chosen arithmetic.
Unobserved state, average reward, undiscounted infinite return,
and general Borel action selection require the additional model
and existence premises appropriate to those different problems.
\end{hbmethod}
